\begin{answers}{ch:rl-llm}{LLM の強化学習：RLHF・DPO・GRPO}
\item \eqref{eq:rl-group-baseline}のアドバンテージは $(0.25,\,0.25,\,0.25,\,-0.75)$、GRPO のアドバンテージは $(0.5,\,0.5,\,0.5,\,-1.5)$ です。
報酬の平均は $\bar r=3/4=0.75$ で、各報酬から引くと前者になります。
標本標準偏差は、平均からのずれの 2 乗の和 $3\times0.25^2+0.75^2=0.75$ を $G-1=3$ で割った $0.25$ の平方根で、$0.5$ です。
前者を $0.5$ で割ると後者になります。
どちらでも、不正解の 1 本の重みの大きさは、正解の 1 本の 3 倍です。
アドバンテージの和は 0 なので、少数派（ここでは 1 本だけの不正解）ほど 1 本あたりの重みが大きくなります。
\item 確率 0.1 のときは約 $2.39$、確率 0.9 のときは約 $0.143$ で、0.1 のほうが罰は約 17 倍大きくなります。
$u=\pi_{\text{ref}}/\pitheta$ なので、0.1 のときは $u=5$、$u-\log u-1=5-1.609-1=2.391$ です。
0.9 のときは $u=0.556$、$0.556-(-0.588)-1=0.143$ です。
参照モデルより確率を下げた方向（$u>1$）では罰が直線的に増えますが、上げた方向（$u<1$）では、$u$ が 0 に近づいても $-\log u$ の増え方がゆるやかだからです（\figref{fig:rl-kl-estimator}）。
\item 数字を 1 つ重ねて書く応答が現れ、本来の正答率は下がりました（実測）。
報酬を \code{1.0 if answer in reply else 0.0} に変え、時間の都合で学習率 $10^{-4}$、40 ステップで GRPO を回しました。
評価用の 200 問では、正しい形「◯◯です。」の割合が 0.870 から 0.625 に下がり、11\% が「正しい数を含むが正解ではない」応答になりました。
例えば「9から4を引くといくつですか？」に「55です。」、「50から3を引くといくつですか？」に「477です。」と答えます。
最後の数字を重ねると、正しい数が応答の中に含まれる見込みが上がるので、緩い報酬はこれに 1 を与えてしまうのです。
同じ学習率で元の報酬を使った本文の実験（\file{out/rl_grpo_highlr.txt}）では、30 ステップで 0.830 だったので、違いは報酬の作り方から来ています。
報酬は「目的そのもの」ではなく、目的が達成できたかを測るための目安にすぎないことがわかります。
\item 縮小した設定で試したところ、どの $\beta$ でも「答えは」で始まる割合は 0 のままで、正答率だけが下がりました（実測）。
\file{reward_model.py} と同じ報酬モデルを学習し、その点数を報酬にして GRPO を 20 ステップ回しました（時間の都合で学習率を本文の 5 倍の $10^{-4}$ にしています）。
正答率は 0.870 から、$\beta=0$ で 0.020、$\beta=0.04$ で 0.030、$\beta=0.5$ で 0.485 になりました。
\file{sft.pt} が「答えは」で書き始める確率はきわめて小さく、サンプリングしてもほぼ出ないので、報酬モデルが最も好む文体には出会えません。
代わりに、報酬モデルが数の違いに付けるわずかな点差（「57です。」が「47です。」より高いなど）を手がかりに、数を書き換えていきました（「12+35は？」に「266です。」）。
報酬の平均は $-2.50$ から $-2.03$ に上がっているので、報酬モデルから見れば「改善」です。
$\beta$ を大きくすると参照モデルから離れにくくなり、正答率の下がり方は遅くなりますが、止まりはしません。
ステップを増やすと $\beta=0.5$ でも下がり続けると予想されます（実測していません）。
\item 数が正しい割合の下がり方は、ずっとゆるやかになります（実測。100 ステップに縮めて実行）。
トークンで見ると、$y_w$ は「答え・は・4・7・です。・\code{<|end|>}」、$y_l$ は「4・7・だよ。・\code{<|end|>}」です。
元の実験の $y_l$「47です。」とは「です。」まで共通でしたが、今度は「だよ。」だけが $y_l$ に固有のトークンです。
DPO は $\log\pitheta(y_l)$ を下げますが、その下げ幅の多くを「だよ。」1 つに負わせられるので、正しい数字を書く確率を削る力が弱まると予想できます。
実際、100 ステップ後の正答率は 0.655 で、元の実験の 100 ステップ目（0.315）より高く保たれました。
\code{rejected} は $-27.7$ と、元の $-8.6$ よりずっと深く下がっており、固有のトークンで差を付けられたことがわかります。
それでも「4」「7」は共通なので、数字を書く力は少し削られ（0.870 から 0.655）、「答えは」で始まる応答も出ませんでした。
\end{answers}
